\documentclass[letterpaper, 10 pt, conference]{ieeeconf}
\IEEEoverridecommandlockouts
\overrideIEEEmargins

\usepackage{cite}
\usepackage{amsmath,amssymb,amsfonts}
\usepackage{algorithmic}
\usepackage{graphicx}
\usepackage{textcomp}
\usepackage{xcolor}
\usepackage{booktabs}
\usepackage{multirow}
\usepackage{subcaption}

\def\BibTeX{{\rm B\kern-.05em{\sc i\kern-.025em b}\kern-.08em
    T\kern-.1667em\lower.7ex\hbox{E}\kern-.125emX}}

\title{Stabilizing Adversarial Self-Play for Pursuit-Evasion Under Field-of-View Constraints}

\author{Anonymous authors}
\begin{document}

\maketitle
\thispagestyle{empty}
\pagestyle{empty}

%%%%%%%%%%%%%%%%%%%%%%%%%%%%%%%%%%%%%%%%%%%%%%%%%%%%%%%%%%%%%%%%%%%%%%%%%%%%%%%%
\begin{abstract}
We study adversarial self-play for 1v1 pursuit-evasion under realistic partial
observability: field-of-view constraints, line-of-sight occlusion, and no external
positioning. A masking curriculum that gradually transitions from full to partial
observability enables stable co-evolution to a fixed point (capture rate
$0.51\pm0.03$ during curriculum decay), but we identify a reproducible
exploitability hole: when the curriculum reaches full partial observability
(\texttt{p\_full}=0), the pursuer policy becomes vulnerable to a best-response
evader (exploitability gap $+0.08$ to $+0.10$ across two independent seeds),
while the evader policy remains robust (best-response pursuer gap negative).
Our key contribution is a set of stabilization mechanisms that extend stable
co-evolution past the curriculum boundary: a visibility reward that teaches the
pursuer to maintain line-of-sight, a forced periodic switching mechanism that
prevents freeze-sticking, and a search staleness reward that encourages
systematic area coverage. Ablations confirm each component is critical; the
best-response analysis characterizes the limits of what curriculum-aided
self-play can guarantee under strict partial observability.
\end{abstract}

%%%%%%%%%%%%%%%%%%%%%%%%%%%%%%%%%%%%%%%%%%%%%%%%%%%%%%%%%%%%%%%%%%%%%%%%%%%%%%%%
\section{Introduction}
\label{sec:introduction}

Pursuit-evasion (PE) games are a fundamental problem in robotics with applications in
surveillance, search-and-rescue, and autonomous navigation. In the 1v1 setting, a pursuer
attempts to capture an evader who tries to escape indefinitely. When both agents are
driven by learned policies trained via self-play, the problem becomes a challenging
adversarial learning problem where training stability is notoriously difficult to achieve.

Most prior work on deep reinforcement learning (RL) for PE assumes full observability:
each agent knows the exact position of the other. This is unrealistic for physical robots
with onboard sensors. Real robots have limited field-of-view (FOV), sensing range, and
occlusion from obstacles. Training under partial observability (PO) introduces new
challenges: the pursuer must learn to search systematically when the evader leaves its
FOV, and the evader must learn to exploit blind spots.

Self-play training under PO is particularly unstable. In our experiments, we found that
standard alternating self-play reliably collapses: one agent dominates, the other receives
vanishing rewards, and co-evolution stalls. We identify two root causes: (1) the pursuer
struggles to bootstrap when it cannot see the evader, and (2) long frozen phases cause
agents to overfit to each other's weaknesses, leading to cyclic collapse.

We propose three stabilization mechanisms:
\begin{enumerate}
    \item \textbf{Visibility reward:} A dense reward that incentivizes the pursuer to maintain
    line-of-sight with the evader, providing learning signal even before capture.
    \item \textbf{Forced periodic switching:} A mechanism that forces role switches every
    2M steps regardless of performance, preventing freeze-sticking and maintaining
    balanced co-evolution.
    \item \textbf{Search staleness reward:} A grid-based coverage reward that encourages
    the pursuer to visit unexplored regions, addressing the critical search skill gap
    under PO.
\end{enumerate}

We evaluate our approach in a simulated 20$\times$20m arena with differential-drive dynamics,
random obstacles, and FOV-based partial observability. Across 5 random seeds, our method
achieves consistent Nash equilibrium convergence (capture rate $0.51 \pm 0.03$) where
prior configurations collapse. Ablations demonstrate that each proposed component is
necessary for stable convergence.

\section{Related Work}
\label{sec:related_work}

\subsection{Pursuit-Evasion with Deep RL}

\subsection{Self-Play and Multi-Agent RL}

\subsection{Partial Observability in RL}

\subsection{Safe RL for Robotics}

\section{Problem Formulation}
\label{sec:problem}

\section{Methods}
\label{sec:methods}

\subsection{Environment}

\subsection{Partial Observability}

\subsection{Self-Play Framework}

\subsection{Stabilization Mechanisms}

\subsubsection{Visibility Reward}

\subsubsection{Forced Periodic Switching}

\subsubsection{Search Staleness Reward}

\subsection{Safety Filter}

\section{Experiments}
\label{sec:experiments}

\subsection{Training Setup}

\subsection{Convergence Results}

\subsection{Ablation Studies}

\subsection{Baseline Comparisons}

\subsection{Behavioral Analysis}

\section{Conclusion}
\label{sec:conclusion}

\section*{Acknowledgments}

\bibliographystyle{IEEEtran}
\bibliography{references}

\end{document}
